\chapter{Depression}

\section*{Definition and Overview}
Depression is a common and serious mental health condition characterised by persistent low mood, loss of interest or pleasure in activities, and a range of emotional, physical, and cognitive symptoms that affect daily functioning. It is more than feeling sad or having a bad day: depression lasts for weeks or longer and interferes with work, study, relationships, and physical health. Depression is a real illness with biological, psychological, and social causes, and it is highly treatable. Most people recover with appropriate treatment, which may include psychological therapy, medication, and lifestyle support. This chapter covers the types, causes, symptoms, and treatment of depression.

\section*{Epidemiology}
Depression is one of the most common mental health conditions in Australia, affecting around one in six adults at some point in their lives, and it is a leading cause of disability worldwide. It is about twice as common in women as in men, and it can begin at any age, though it most often starts in adolescence or early adulthood. Around one in ten new mothers experiences depression after childbirth. Depression frequently coexists with anxiety, substance use, and chronic physical illness, and it increases the risk of suicide. Despite being common and treatable, many people do not seek help, often because of stigma or difficulty accessing services.

\section*{Aetiology and Pathophysiology}
\begin{itemize}
    \item \textbf{Genetic factors:} depression runs in families, and many genes contribute to risk
    \item \textbf{Brain chemistry:} imbalances in neurotransmitters such as serotonin, noradrenaline, and dopamine are involved
    \item \textbf{Hormonal changes:} thyroid problems, pregnancy, childbirth, and menopause can trigger depression
    \item \textbf{Stressful life events:} loss, relationship breakdown, financial stress, and trauma are common triggers
    \item \textbf{Physical illness:} chronic pain, heart disease, cancer, and other conditions increase risk
    \item \textbf{Medications and substances:} some medicines, alcohol, and drugs can cause or worsen depression
    \item \textbf{Personality and coping:} low self-esteem, perfectionism, and poor coping skills increase vulnerability
    \item \textbf{Social factors:} isolation, unemployment, and disadvantage contribute
\end{itemize}

\section*{Clinical Features}
\begin{itemize}
    \item Persistent low, sad, or empty mood lasting most of the day, nearly every day
    \item Loss of interest or pleasure in activities that were once enjoyable
    \item Significant changes in appetite and weight
    \item Sleep disturbance: insomnia, especially early-morning waking, or sleeping too much
    \item Fatigue and loss of energy
    \item Difficulty concentrating, making decisions, or remembering
    \item Feelings of worthlessness, excessive guilt, or hopelessness
    \item Slowed movement or agitation
    \item Withdrawal from family, friends, and activities
    \item Physical symptoms such as headaches, aches, and digestive problems
    \item Recurrent thoughts of death or suicide
\end{itemize}

\section*{Differential Diagnosis}
\begin{itemize}
    \item Normal sadness and grief, which are less persistent and less impairing
    \item Bipolar disorder, in which depressive episodes alternate with mania or hypomania
    \item Anxiety disorders, which commonly coexist with depression
    \item Medical conditions causing similar symptoms, such as thyroid disease, anaemia, and vitamin B12 deficiency
    \item Medication side effects and substance use
    \item Adjustment disorder, in which distress follows a stressor but is less severe
    \item Premenstrual dysphoric disorder and postnatal depression
\end{itemize}

\section*{When to Seek Medical Care}
Anyone who has had low mood, loss of interest, or other symptoms of depression for more than two weeks, or whose symptoms affect work, study, or relationships, should see a GP. A GP can assess symptoms, arrange referral to a psychologist or psychiatrist, and prescribe treatment. Seeking help early improves outcomes.

\textbf{Call 000 (or local emergency services) for:}
\begin{itemize}
    \item Thoughts of self-harm or suicide, or concerns that someone is in danger
    \item Severe agitation, confusion, or inability to care for oneself
    \item Any mental health crisis in which a person is at immediate risk
\end{itemize}
\textbf{Support lines:} Lifeline (13 11 14), Beyond Blue (1300 224 636), and Suicide Call Back Service (1300 659 467) are available 24 hours.

\section*{Diagnosis and Investigations}
\begin{itemize}
    \item \textbf{Clinical interview:} a detailed discussion of mood, thoughts, sleep, appetite, energy, and functioning
    \item \textbf{Screening tools:} questionnaires such as the PHQ-9 help assess severity and monitor response
    \item \textbf{Physical examination and blood tests:} thyroid function, full blood count, and vitamin levels exclude medical causes
    \item \textbf{Risk assessment:} careful enquiry about suicidal thoughts and self-harm
    \item \textbf{Review of medications and substances:} identifying contributing factors
    \item \textbf{Referral:} to a psychologist or psychiatrist for specialised assessment and treatment
\end{itemize}

\section*{Management}
\begin{itemize}
    \item \textbf{Psychological therapy:} cognitive-behavioural therapy (CBT), interpersonal therapy, and other evidence-based therapies are highly effective
    \item \textbf{Antidepressant medication:} effective for moderate to severe depression; usually takes several weeks to work and should be continued for months after recovery
    \item \textbf{GP Mental Health Treatment Plan:} provides Medicare-subsidised sessions with a psychologist
    \item \textbf{Lifestyle measures:} regular exercise, good sleep, healthy eating, and reduced alcohol
    \item \textbf{Social support:} staying connected with family, friends, and support groups
    \item \textbf{Managing contributing factors:} treatment of physical illness, insomnia, or substance use
    \item \textbf{Hospital care:} for severe depression, high risk, or inability to care for oneself
    \item \textbf{Specialist treatments:} for treatment-resistant depression, including different medications, augmentation, and in some cases electroconvulsive therapy (ECT) or other procedures
\end{itemize}

\section*{Complications}
\begin{itemize}
    \item Suicide and self-harm, the most serious complications
    \item Impaired work performance, job loss, and financial difficulty
    \item Relationship breakdown and social isolation
    \item Alcohol and drug use as self-medication
    \item Physical health problems, including heart disease, diabetes, and weakened immunity
    \item Chronic or recurrent depression and treatment resistance
    \item Effects on children and families of a parent with depression
\end{itemize}

\section*{Prognosis and Outcomes}
With appropriate treatment, most people with depression recover. Psychological therapy and medication are each effective, and combining them is often best for moderate to severe depression. Around 50--70\% of people respond to first-line treatment, and many achieve full remission. Depression can recur, so continuing treatment after recovery and learning relapse-prevention strategies are important. With early help and good support, most people return to their previous level of functioning and live full, productive lives.

\section*{Prevention and Self-Care}
\begin{itemize}
    \item Maintain regular exercise, which is protective against depression and improves mood
    \item Prioritise sleep and manage stress with relaxation or mindfulness
    \item Stay connected with supportive people and engage in meaningful activities
    \item Limit alcohol and avoid recreational drugs
    \item Seek help early when symptoms appear rather than waiting
    \item Take prescribed treatment consistently and attend follow-up appointments
    \item Learn to recognise early warning signs of relapse
    \item Support others by listening without judgement and encouraging professional help
\end{itemize}

\section*{Sources and Further Reading}
The following reputable sources provide further information for healthcare students and patients seeking reliable, current advice about depression.

\begin{itemize}
    \item Beyond Blue. \textit{Depression: symptoms, treatment, and support.} https://www.beyondblue.org.au/
    \item Healthdirect Australia. \textit{Depression: information and support.} https://www.healthdirect.gov.au/depression
    \item Black Dog Institute. \textit{Depression and mood disorders: information and help.} https://www.blackdoginstitute.org.au/
    \item NHS. \textit{Clinical depression: symptoms and treatment.} https://www.nhs.uk/mental-health/conditions/clinical-depression/
    \item Mayo Clinic. \textit{Depression: diagnosis and treatment.} https://www.mayoclinic.org/diseases-conditions/depression/
    \item MSD Manual. \textit{Depressive disorders: professional overview.} https://www.msdmanuals.com/
\end{itemize}
